% !TEX root = ../main.tex
\section{Introduction}
Mathematical reasoning is a demanding test of non-autoregressive generation. A solution must maintain dependencies across intermediate calculations, and a locally plausible continuation can still produce an incorrect final answer. Autoregressive models express these dependencies through sequential token generation. Diffusion models offer a different computation pattern: they repeatedly refine a representation of the whole answer, allowing the number of refinement steps to vary independently of its token length.

Much of the progress in diffusion language modeling uses discrete masking, including LLaDA \citep{llada}. Continuous alternatives transport noisy vectors toward language representations through a learned flow. Recent work has explored hyperspherical embeddings, masked continuous paths, and flow-map refinement \citep{hyperspherical,mlfm,posterior}. ELF \citep{elf} provides a particularly direct formulation: most computation occurs in a continuous contextual embedding space, followed by a shared-backbone token decoder. This raises a practical question: how far can a relatively simple continuous denoising recipe go on verifiable reasoning tasks?

We study three coupled choices. First, the representation must retain information useful for predicting the answer, rather than only directions with high variance. We learn low-dimensional projectors by matching a frozen teacher's downstream token distributions and explicitly whiten their coordinates. We combine layers 16, 24, and 32 of Qwen3-4B-Instruct-2507, allocating a total of 1,024 latent dimensions per token. The choice is guided by a pronounced representation transition after layer 16 and regular spacing through deeper layers.

Second, conditioning must remain effective when the teacher is removed from the inference path. We replace teacher prompt activations with a six-layer prompt Transformer. Learning the prompt representation is staged: imitation supplies an initialization, after which the prompt encoder and an already-trained denoiser adapt jointly to the reasoning objective. This separates learning to generate answers from learning a new conditioning interface.

Third, inference compute can be spent either on refining one answer or on drawing additional answers. We evaluate both oracle coverage and final-answer voting because a correct answer appearing somewhere in a sample set does not by itself specify which answer to return. We also examine how reward-based fine-tuning changes this tradeoff.

Our contributions are: (i) a continuous latent-diffusion recipe for mathematical reasoning that combines teacher-trained projections, covariance whitening, and staged prompt adaptation; (ii) controlled measurements of prompt substitution, early joint training, layer-clock schedules, and reward-based fine-tuning; and (iii) a seeded evaluation of denoising depth versus sample count, including answer selection, diversity, and measured inference cost. Our large denoiser reaches 63.74\% mean pass@1 on GSM8K and 23.60\% on MATH500 at 64 steps. The inference model uses a learned prompt encoder and a frozen token lookup table, without running the teacher's Transformer layers.
